% SEGMENT OLP-0051-S01
% Part: sets-functions-relations
% Chapter: infinite
% Section: dedekind
%
\documentclass[../../../include/open-logic-section]{subfiles}

\begin{document}

\olfileid{sfr}{infinite}{dedekind}
\olsection{டெடெகிண்ட் இயற்கணிதங்கள்}

இயல் எண்கள் முடிவுறாதவையாக மட்டும் இருக்க வேண்டும் என்று விரும்பவில்லை;
அவை சில (இயற்கணித) பண்புகளையும் கொண்டிருக்க வேண்டும் என்று விரும்புகிறோம்:
கூட்டல், பெருக்கல் முதலியவற்றின் கீழ் அவை ஒழுங்காகச் செயல்பட வேண்டும்.

\emph{தொடரிச் சார்பு} என்ற கருத்தை அடிப்படையாக எடுத்துக்கொண்டு,
பின்வரும் பண்புகளைக் கொண்டவையாக எண்களைப் பண்புரைப்பதே
டெடெகிண்டின் கருத்து:
\begin{enumerate}
	\item எந்த எண்ணின் தொடரியும் அல்லாத $0$ என்ற ஓர் எண் உள்ளது
	\\ அதாவது, $0 \notin \ran{s}$
	\\ அதாவது, $\forall x\ s(x) \neq 0$
	\item வெவ்வேறு எண்களுக்கு வெவ்வேறு தொடரிகள் உள்ளன
	\\ அதாவது, $s$ ஓர் !!a{injection}
	\\ அதாவது, $\forall x \forall y (s(x) = s(y) \lif x = y)$
	\item\ollabel{repeatedapplication} தொடரிச் சார்பை $0$ மீது
    மீண்டும் மீண்டும் பயன்படுத்துவதன் மூலம் ஒவ்வோர் எண்ணும் பெறப்படுகிறது.
\end{enumerate}
முதலிரண்டு நிபந்தனைகளையும் முதல்தரத் தருக்கவியலைப் பயன்படுத்தி
எளிதாகக் கையாளலாம் (மேலே பார்க்கவும்).
ஆனால் முதல்தரத் தருக்கவியலை மட்டும் கொண்டு
\olref{repeatedapplication} ஐக் கையாள முடியாது.
\olref{repeatedapplication} என்ற நிபந்தனையைப் பின்வருமாறு
கணக்கோட்பாட்டு முறையில் மறுவடிவமைத்ததே டெடெகிண்டின் முன்னேற்றம்:
\begin{enumerate}
	\item[3$'$.] இயல் எண்கள், தொடரிச் சார்பின் கீழ்
	\emph{மூடப்பட்டுள்ள} மீச்சிறு கணம்;
    அதாவது, அந்தக் கணத்தின் எந்த !!{element} மீதும் $s$ ஐப்
    பயன்படுத்தினால், அந்தக் கணத்தின் மற்றோர் !!{element}[acc]ப் பெறுகிறோம்.
\end{enumerate}
ஆனால் இதனை மெதுவாக விரித்துரைக்க வேண்டும்.

% SEGMENT OLP-0051-S02
\begin{defn}\ollabel{Closure}
	எந்தச் சார்பு $f$ க்கும், $(\forall x \in X)f(x) \in X$ என
    இருக்கும்போது, அப்போதும் அப்போது மட்டுமே, கணம் $X$ ஆனது
    $f$-\emph{மூடப்பட்டது} ஆகும்.
    இப்போது எந்த $o$ க்கும் பின்வருமாறு வரையறுக்க:
	$$\closureofunder{f}{o} = \bigcap\Setabs{X}{o \in X\text{ மற்றும் }X
\text{ ஆனது $f$-மூடப்பட்டது}}$$
\end{defn}

% SEGMENT OLP-0051-S03
எனவே $\closureofunder{f}{o}$ என்பது, $o$ ஐ ஓர் !!a{element} ஆகக்
கொண்ட எல்லா $f$-மூடிய கணங்களின் வெட்டாகும்.
அப்படியானால், உள்ளுணர்வின்படி, $\closureofunder{f}{o}$ என்பது
$o$ ஐ ஓர் !!a{element} ஆகக் கொண்ட \emph{மீச்சிறு} $f$-மூடிய கணம்.
பின்வரும் முடிவு அந்த உள்ளுணர்வைத் துல்லியமாக்குகிறது:
\begin{lem}\ollabel{closureproperties}
	எந்தச் சார்பு $f$ க்கும், அதன் சார்பகமும் வீச்சகமும் ஒரே கணம்
    $A$ ஆகவும், எந்த $o \in A$ க்கும்:
	\begin{enumerate}
		\item\ollabel{closurehaselem} $o \in \closureofunder{f}{o}$; மேலும்
		\item\ollabel{closureclosed} $\closureofunder{f}{o}$ ஆனது $f$-மூடப்பட்டது; மேலும்
		\item\ollabel{closuresmallest} $X$ ஆனது $f$-மூடப்பட்டும் $o \in
		X$ எனவும் இருந்தால், $\closureofunder{f}{o} \subseteq X$
	\end{enumerate}
\end{lem}

% SEGMENT OLP-0051-S04
\begin{proof}
$o$ ஐ ஓர் !!a{element} ஆகக் கொண்ட குறைந்தது ஒரு $f$-மூடிய கணம்
உள்ளது என்பதை கவனிக்கவும்; அது $\ran{f}\cup \{o\}$.
எனவே இத்தகைய கணங்கள் \emph{அனைத்தின்} வெட்டான
$\closureofunder{f}{o}$ உள்ளது.
இப்போது \olref{closurehaselem}--\olref{closuresmallest} ஐச்
சரிபார்க்க வேண்டும்.

\olref{closurehaselem} ஐப் பொறுத்தவரை:
$o$ ஐ ஓர் !!a{element} ஆகக் கொண்ட கணங்களின் வெட்டாக இருப்பதால்,
$o \in \closureofunder{f}{o}$.

\olref{closureclosed} ஐப் பொறுத்தவரை:
$x \in \closureofunder{f}{o}$ எனக் கொள்க.
எனவே $o \in X$ என்றும் $X$ ஆனது $f$-மூடப்பட்டது என்றும் இருந்தால்,
$x \in X$; மேலும் $X$ ஆனது $f$-மூடப்பட்டது என்பதால் $f(x) \in X$.
ஆகவே $f(x) \in \closureofunder{f}{o}$.

\olref{closuresmallest} ஐப் பொறுத்தவரை:
பொதுவாகவே, $X \in C$ எனில் $\bigcap C \subseteq X$.
\end{proof}

% SEGMENT OLP-0051-S05
இதனைப் பயன்படுத்தி பின்வருமாறு கூறலாம்:

\begin{defn}
ஒரு \emph{டெடெகிண்ட் இயற்கணிதம்} என்பது,
$f \colon A \to A$ என்ற ஒரு சார்பையும் ஏதோ ஓர் $o \in A$ ஐயும்
கொண்ட ஒரு கணம் $A$; அது பின்வருவனவற்றை நிறைவேற்றுகிறது:
	\begin{enumerate}
		\item \ollabel{ded:proper} $o \notin \ran{f}$
		\item \ollabel{ded:injection} $f$ ஓர் !!a{injection}
		\item \ollabel{ded:closure} $A = \closureofunder{f}{o}$
	\end{enumerate}
\end{defn}

% SEGMENT OLP-0051-S06
$A = \closureofunder{f}{o}$ என்பதால், $A$ என்பது $o$ ஐ ஓர்
!!a{element} ஆகக் கொண்ட மீச்சிறு $f$-மூடிய கணம் என நமது முந்தைய
முடிவு கூறுகிறது.
ஒரு டெடெகிண்ட் இயற்கணிதம் டெடெகிண்ட்-முடிவுறாதது என்பது தெளிவு;
வரையறையின் \olref{ded:proper}, \olref{ded:injection} என்ற
கூறுகளை மட்டும் பார்க்கவும்.
ஆனால் எந்த டெடெகிண்ட்-முடிவுறாத கணத்தையும் ஒரு டெடெகிண்ட்
இயற்கணிதமாக மாற்றலாம் என்பது மேலும் சுவையான உண்மை.

% SEGMENT OLP-0051-S07
\begin{thm}\ollabel{thm:DedekindInfiniteAlgebra}
ஒரு டெடெகிண்ட்-முடிவுறாத கணம் இருக்குமானால்,
ஒரு டெடெகிண்ட் இயற்கணிதம் உள்ளது.
\end{thm}

% SEGMENT OLP-0051-S08
\begin{proof}
$D$ டெடெகிண்ட்-முடிவுறாததாக இருக்கட்டும்.
எனவே $g \colon D \to D$ என்ற ஓர் உட்செலுத்தலும்
$o  \in D \setminus \ran{g}$ என்ற ஓர் உறுப்பும் உள்ளன.
இப்போது $A = \closureofunder{g}{o}$ என்க;
\olref{closureproperties} இன்படி $A$ உள்ளது, $o \in A$.
$f = \funrestrictionto{g}{A}$ என்க.
$A, f, o$ ஆகியவை ஒரு டெடெகிண்ட் இயற்கணிதத்தை அமைக்கின்றன
எனக் காட்டுவோம்.

\olref{ded:proper} ஐப் பொறுத்தவரை:
$o \notin \ran{g}$, $\ran{f} \subseteq \ran{g}$;
எனவே $o\notin \ran{f}$.

\olref{ded:injection} ஐப் பொறுத்தவரை:
$g$ என்பது $D$ மீதான ஓர் உட்செலுத்தல்;
எனவே $f \subseteq g$ உம் ஓர் உட்செலுத்தலாக இருக்க வேண்டும்.

\olref{ded:closure} ஐப் பொறுத்தவரை:
\olref{closureproperties} இன்படி $A$ ஆனது $g$-மூடப்பட்டது;
எனவே இன்னும் வலுவாக, $A$ ஆனது $f$-மூடப்பட்டது.
ஆகவே \olref{closureproperties} இன்படி
$\closureofunder{f}{o} \subseteq A$.
மேலும் $\closureofunder{f}{o}$ ஆனது $f$-மூடப்பட்டது;
$f = \funrestrictionto{g}{A}$ என்பதால்,
$\closureofunder{f}{o}$ ஆனது $g$-மூடப்பட்டது எனப் பின்தொடரும்.
எனவே \olref{closureproperties} இன்படி
$A \subseteq \closureofunder{f}{o}$.
\end{proof}

\end{document}
